\documentclass[11pt]{article}
\usepackage[margin=1in]{geometry}
\usepackage{amsmath,amssymb,amsthm}
\usepackage{hyperref}
\newtheorem{theorem}{Theorem}
\newtheorem{lemma}[theorem]{Lemma}
\newtheorem{proposition}[theorem]{Proposition}
\newtheorem{remark}[theorem]{Remark}
\newcommand{\an}{\mathrm{an}}
\newcommand{\Int}{\operatorname{Int}}
\newcommand{\Fin}{\operatorname{Fin}}
\newcommand{\dd}{\mathrm d'\mathrm d''}
\title{A relative-boundary obstruction to sequential classical-psh regularization\\
\large Partial progress for P02.1--P02.4}
\author{Research draft}
\date{3 October 2026}
\begin{document}
\maketitle
\begin{abstract}
Over $K=k((\pi))$, where $k$ is an uncountable algebraically closed
coefficient field and $|\pi|=e^{-1}$, we give a finite upper semicontinuous
classically plurisubharmonic function on a good strict analytic polydisc
of any dimension at least two which, near a specified relative-boundary
point, is not the pointwise limit of any sequence of continuous real
functions. Thus decreasing sequences of semipositive real PL, finite
locally Fubini--Study, weakly smooth psh, and ordinary CLD-smooth psh
functions all fail for this example. Separate explicit decreasing nets
of locally Fubini--Study semipositive PL functions and of ordinary
CLD-smooth psh functions do converge to it. When $k$ is countable, the
same nets admit decreasing cofinal sequences. The result is a partial
obstruction to sequential formulations, not a characterization of spaces;
it does not refute formulations allowing arbitrary nets. None of the
four complete atomic targets is claimed solved.
\end{abstract}

\section{Conventions and imported statements}
We use Berkovich analytic spaces, their ordinary topology, and
$v=-\log|\cdot|$. A relative boundary is always relative to the ground
field, and need not be a topological boundary in the space itself.
Classical psh has exactly the meaning of \cite[Definition~4.1]{GR}:
an upper semicontinuous function $h:Y\to\mathbb R\cup\{-\infty\}$ is
tested after every analytic extension $L/K$ and every morphism from an
$L$-analytic curve to $Y_L$. Its pullback must be subharmonic in the
functorial sense of \cite[Definition~3.4]{GR}. That definition in turn
tests smooth connected strict curves over nontrivially valued analytic
extensions, allowing the identically $-\infty$ function.

A finite locally Fubini--Study function means the prescribed local
finite logarithmic maximum of analytic functions, with a positive
integer denominator and no common zero. We only construct expressions
of denominator $1$. Real PL functions use the real extension of the
local integral PL class. Semipositivity is residue-line-bundle nefness.
Ordinary CLD-smooth functions use analytic moment maps to split tori;
weakly smooth functions permit harmonic tropicalizations. We only need
their respective definitional continuity for the negative result.
Ordinary smooth psh in the constructive result is certified by positivity
of its smooth $(1,1)$-form $\dd h$. No identification of the two smooth
classes or of different positivity notions is assumed.

Here are the essential external facts, with their exact scopes. They
are imported from a primary source, not independently established by
an audit receipt or by the catalogue description.
\begin{enumerate}
\item The topology of $\mathcal M(A)$ is generated by the functions
$|f|$, $f\in A$, and its Shilov boundary is the unique minimal set on
which every analytic function attains its maximum
\cite[\S2.4]{GR}. Distinct seminorms are separated by an evaluation,
so this topology is Hausdorff.
\item An analytic curve means a separated pure one-dimensional analytic
space. For an affinoid analytic curve its relative boundary equals its
Shilov boundary \cite[\S3.1(3)]{GR}.
\item If $C$ is an analytic curve and $h:C\to\mathbb R\cup\{-\infty\}$
is upper semicontinuous and subharmonic on $\Int(C)$, then $h$ is
subharmonic on $C$ \cite[Proposition~3.5(3)]{GR}. This conclusion uses
Definition~3.4, including the scalar extensions, rather than just a
nonnegative Laplacian on the original interior.
\item On an analytic curve, classical psh is equivalent to subharmonicity
\cite[\S4.1.2]{GR}. For arbitrary analytic spaces, logarithms of absolute
values of analytic functions are classically psh
\cite[Proposition~4.2(2)]{GR}; nonnegative linear combinations and finite
maxima preserve classical psh \cite[Proposition~4.2(3)]{GR}.
Classical psh is preserved by analytic pullback and is equivalent to
classical psh after any analytic extension
\cite[Proposition~4.2(5),(6), forward direction of (6)]{GR}.
\item If $K$ is nontrivially valued, $Y$ is good and strictly $K$-analytic,
and $h:Y\to\mathbb R$ is real PL, then $h$ is semipositive if and only
if it is classically psh \cite[Theorem~4.3]{GR}.
\item For $Y\xrightarrow{f}Z\to S$, one has
$\partial(Y/S)\subset\partial(Y/Z)\cup f^{-1}(\partial(Z/S))$,
with equality when $Z\to S$ is locally separated
\cite[\S2.11(4), equations (2.11.3)--(2.11.4)]{GR}.
\end{enumerate}
The strict closed polydisc
$\mathcal M(K\langle T_1,\ldots,T_n\rangle)$ is good, strict, separated
and pure dimension $n$, by the standard Tate-algebra geometry
\cite[Chapter~2, \S2.1]{Ber}. We use precisely this polydisc. In
particular, no boundarylessness is implicit in its regularity.

\section{The analytic example}
Let $k$ be algebraically closed with its trivial absolute value, and put
$K=k((\pi))$ with $|\pi|=e^{-1}$. This is complete and nontrivially
discretely valued. Write
\[
D=\mathcal M(K\langle T\rangle),\qquad
X=\mathcal M(K\langle T_1,\ldots,T_n\rangle),\quad n\ge2.
\]
Let $\eta\in D$ be the Gauss point, let $p:X\to D$ be first-coordinate
projection, and let $i:D\to X$ be the closed analytic embedding
$T_1=T$, $T_2=\cdots=T_n=0$. Set
\[
w(y)=\begin{cases}1&y=\eta,\\0&y\ne\eta,\end{cases}
\qquad u=w\circ p,\qquad x=i(\eta).
\]

\begin{theorem}\label{thm:obstruction}
Suppose $k$ is uncountable. The function $u$ is finite, bounded, upper
semicontinuous and classically psh, and $x\in\partial X$. For every
ordinary open neighborhood $W$ of $x$, no sequence of continuous real
functions on $W$ converges pointwise to $u|_W$. Consequently none of
P02.1--P02.4 can hold with decreasing sequences on every good strict
space when these fields and relative-boundary points are allowed.
This is a failure example, not the full characterization requested by
any of those targets.
\end{theorem}

\begin{lemma}\label{lem:cofinite}
Every ordinary neighborhood of $\eta$ in $D$ contains all but finitely
many of the $K$-rational coefficient points $T=a$, $a\in k$.
\end{lemma}
\begin{proof}
For nonzero $f=\sum c_jT^j\in K\langle T\rangle$, its Gauss norm
$\|f\|=\max_j|c_j|$ is attained because $c_j\to0$, and lies in
$|K^\times|$. Normalize $f$ to norm $1$. Its reduction $\bar f\in k[T]$
is a nonzero polynomial. For $a$ outside its finite root set,
$|f(a)|=1=|f(\eta)|$. Undoing the normalization preserves the equality.
The zero function imposes no exceptions. A basic ordinary neighborhood
of $\eta$ consists of conditions on finitely many evaluations. Outside
the union of their finite exceptional sets all evaluations equal their
values at $\eta$, so all conditions hold. These basic sets generate the
ordinary topology.
\end{proof}

\begin{proof}[Proof of Theorem~\ref{thm:obstruction}]
The Gauss norm is multiplicative: after normalization, reduction turns
the product of two nonzero norm-one Tate series into a product of nonzero
polynomials in the integral domain $k[T]$. Every point of $D$ satisfies
$|f(y)|\le\|f\|=|f(\eta)|$. Hence $\{\eta\}$ is a Shilov boundary,
and minimality gives that it is the Shilov boundary. The imported
curve-boundary fact gives $\partial D=\{\eta\}$.

The singleton $\{\eta\}$ is closed. Thus $w$ is upper semicontinuous,
and $w=0$ on $\Int(D)$. Constants are subharmonic. Applying exactly
\cite[Proposition~3.5(3)]{GR} gives subharmonicity on all of $D$, and
\cite[\S4.1.2]{GR} gives classical psh. Pullback by $p$, using
\cite[Proposition~4.2(6)]{GR}, proves that $u$ is classically psh. In
particular, all scalar-extension curve tests are included. Its finite
values and upper semicontinuity also follow directly from its definition.
Since the structural morphism of $D$ is separated, the boundary formula
in \cite[\S2.11(4)]{GR} implies $p^{-1}(\partial D)\subset\partial X$,
so $x\in\partial X$.

Now suppose $h_j\in C(W,\mathbb R)$ converge pointwise to $u|_W$.
The inverse image $U=i^{-1}(W)$ is an ordinary neighborhood of $\eta$.
Continuity at $x$ gives, for each $j$, a neighborhood $N_j\subset U$
of $\eta$ such that
\[
|h_j(i(y))-h_j(x)|<\tfrac14\quad(y\in N_j).
\]
Lemma~\ref{lem:cofinite} gives a finite exceptional set $E_j\subset k$
outside which every coefficient point belongs to $N_j$. Uncountability
allows $a\in k\setminus\bigcup_jE_j$. Then $i(a)\in W$ and
$u(i(a))=0$, since $a\ne\eta$: the function $T-a$ has norm $0$ at
$a$ and norm $1$ at $\eta$. But $h_j(x)\to1$, so the displayed estimate
forces $h_j(i(a))>1/2$ eventually, contrary to $h_j(i(a))\to0$.
Monotonicity was not used.

Real PL functions and finite local logarithmic maxima are continuous.
Ordinary and weak smoothness each imply continuity from their respective
local definitions. Thus all four decreasing sequential conclusions fail.
Even if a broader logarithmic-max class permits a base locus, a
decreasing approximation to $u$ must have each term finite at $x$ and
therefore continuous on a neighborhood of $x$. The preceding argument
only needs such individual neighborhoods, so also rules out that case.
\end{proof}

\section{Explicit nets: an essential limitation of the obstruction}
We now assume only that $k$ is algebraically closed, without a
cardinality restriction.

\begin{lemma}\label{lem:detection}
For every $y\ne\eta$ in $D$ there is $a\in k$ with $|T-a|(y)<1$.
\end{lemma}
\begin{proof}
Distinct seminorms give $f\in K\langle T\rangle$ with
$|f(y)|<\|f\|$; no point can exceed the Gauss norm. Normalize to
$\|f\|=1$. Lift its nonzero reduction to $q\in k[T]$, so
$\|f-q\|\le|\pi|<1$. Thus $|q(y)|<1$. Factor
$q=c\prod_a(T-a)^{e_a}$ over the algebraically closed coefficient field.
Here $|c|=1$ and all $|T-a|(y)\le1$. Some factor must have norm $<1$;
if $q$ were constant it would have norm $1$, already a contradiction.
\end{proof}

\begin{proposition}\label{prop:nets}
Both $w$ on $D$ and $u$ on $X$ are pointwise decreasing limits of a
net of finite locally Fubini--Study semipositive real-PL functions,
and, separately, a net of finite ordinary CLD-smooth psh functions.
The formulas persist under every complete scalar extension, and the
extended nets converge to the pullbacks of $w$ and $u$.
\end{proposition}
\begin{proof}
Direct $I=\Fin(k)\times\mathbb N_{\ge1}$ by inclusion and the usual
integer order. For finite $F\subset k$, let
\[
P_F=\prod_{a\in F}(T-a),\quad P_\varnothing=1,\qquad
v_{F,m}=\max\{0,1+m\log|P_F|\}
=\log\max\{1,|\pi^{-1}P_F^m|\}.
\]
At zeros the value is $0$. The last expression is a finite logarithmic
maximum of analytic sections $1$ and $\pi^{-1}P_F^m$, denominator $1$,
with no common zero. Hence it is locally Fubini--Study. It is real PL
on the strict Laurent cover given by $|\pi^{-1}P_F^m|\le1$ and
$|\pi^{-1}P_F^m|\ge1$: the expressions are respectively $0$ and
$\log|\pi^{-1}P_F^m|$, an analytic-unit logarithm on the second domain,
and agree on the overlap. The same expressions in $T_1$ work on $X$.
They are classically psh by \cite[Proposition~4.2(2),(3)]{GR}, so they
are semipositive by \cite[Theorem~4.3]{GR}. Nontriviality, goodness,
strictness and the real-PL premise of that theorem have all been checked.

Every $|T-a|$ is at most $1$ on $D$. Increasing $F$ or $m$ decreases
$1+m\log|P_F|$ and hence decreases $v_{F,m}$. At $\eta$ the value is
$1$. At any other point, Lemma~\ref{lem:detection} provides a factor of
norm $<1$; once it is included, increasing $m$ forces the argument to
$-\infty$ or it is already $-\infty$. The value is eventually $0$.
Thus $v_{F,m}\downarrow w$ and $v_{F,m}\circ p\downarrow u$.

For ordinary smooth approximants, choose a nonnegative even smooth
probability density $\rho$ supported in $[-1/4,1/4]$, for instance the
normalized function $\exp(-1/(1-16s^2))$ on $|s|<1/4$, extended by
zero. Define
\[
\lambda(t)=\int_{\mathbb R}\max\{0,t+s\}\rho(s)\,ds.
\]
Then
\[
\lambda'(t)=\int_{s>-t}\rho(s)\,ds\ge0,\qquad
\lambda''(t)=\rho(-t)\ge0.
\]
In particular $\lambda$ is smooth, nondecreasing and convex. It is
$0$ for $t\le-1/4$ and equals $t$ for $t\ge1/4$, because $\rho$ has
mean zero. Put
\[
q_{F,m}=\lambda(1+m\log|P_F|),
\]
with value $0$ at zeros. Near a zero of $P_F$, the formula is identically
$0$ by the negative constant tail. On the nonzero locus $P_F$ is a
unit on a neighborhood; this follows by taking a Laurent neighborhood
with a positive lower bound on $|P_F|$. There $q_{F,m}$ is the pullback
of the ordinary smooth convex function $t\mapsto\lambda(1-mt)$ through
the analytic moment map $t=-\log|P_F|$. Its second derivative is
$m^2\lambda''(1-mt)\ge0$. The CLD definition therefore gives ordinary
smoothness and positive smooth Hessian. Equivalently, on the unit locus,
\[
\dd q_{F,m}
=m^2\lambda''(1+m\log|P_F|)
\,\mathrm d'\log|P_F|\wedge\mathrm d''\log|P_F|\ge0.
\]
At the zeros the form is locally zero. These local descriptions agree
with one global expression. The same construction in $T_1$ certifies
ordinary smooth psh on $X$, including its relative boundary, without
pulling back any current.

Monotonicity follows from the nondecreasing $\lambda$. At $\eta$ the
value is $\lambda(1)=1$, and at every other point it is eventually $0$
by the detection lemma and the negative tail. This proves convergence.

After any complete extension $L/K$, the identical analytic formulas
make sense on $D_L$ and $X_L$. A norm of a $K$-defined function at a
point upstairs equals its norm at its image downstairs. Hence the
displayed nets converge to the pulled-back functions. The unit-locus,
constant-tail, and smooth-Hessian checks are unchanged; classical psh
and the good-strict real-PL comparison also apply over $L$. No equation
of arbitrary currents has been transported through a morphism.
\end{proof}

\begin{proposition}\label{prop:threshold}
For this construction, when $k$ is countable, the two nets in
Proposition~\ref{prop:nets} have decreasing cofinal sequences. When $k$
is uncountable, $I$ has no countable cofinal subset, and no sequence of
continuous functions unrelated to the nets can approximate $u$ near $x$.
\end{proposition}
\begin{proof}
If $k=\{a_1,a_2,\ldots\}$, take $F_j=\{a_1,\ldots,a_j\}$ and $m_j=j$.
For every finite $F$ and every $m$, eventually $F\subset F_j$ and
$m\le j$. This gives the asserted cofinal sequences and convergence.
If $k$ is uncountable, the union of a countable family of finite $F_j$
is countable and omits some $a$. Such a family cannot dominate the
index $(\{a\},1)$. The stronger impossibility is
Theorem~\ref{thm:obstruction}.
\end{proof}

\section{Comparison, limitations and verification}
The supplied family dossier asks for geometric characterizations for
arbitrary classical psh functions. Its toric endpoint theorems start
with finite continuous invariant current-psh functions on ordinary open
torus domains. The weakly smooth curve theorem of Gubler--Jell--Rabinoff,
Proposition~15.11, treats arbitrary Thuillier-subharmonic functions on
smooth boundaryless curves and gives a decreasing net; neither continuity
nor a proper algebraic curve is required. A separate ordinary-smooth
curve endpoint concerns continuous current-psh functions.
Our example is a discontinuous boundary jump on a strict affinoid
polydisc. It does not contradict those known cases and is not counted
as one of them. No novelty or priority claim over the entire literature
is made.

The four sequential failures use only continuity, so no ordinary/weak
smoothness identification is involved. For weak smoothness we need no
constructive net and no comparison of its positivity notion with CLD
positivity. The nets actually constructed certify semipositive PL,
locally Fubini--Study, and ordinary smooth psh separately. None converges
locally uniformly at $x$, since $w$ is discontinuous and every
neighborhood of $\eta$ contains coefficient points.

This is an admissible analytic counterexample to a universal sequential
claim allowing these boundaries and fields. It is \emph{partial}
progress for each atom: the space characterization, the boundaryless
higher-dimensional problem, arbitrary continuous noninvariant starting
functions, and the general net formulation all remain unresolved.
The sharp cardinality comparison is only for the displayed Laurent-field
boundary jump, not a sharp classification of all failure ranges.

The primary PDF used is arXiv:2506.13548v2, with SHA-256
\begin{center}\small\ttfamily
83ca9a19394f44bb2e4e1aa316de95938f31\\
1acd53a00fe250fa0f605184b726
\end{center}
Its statements were checked in the local primary text. The independent
plain-text full argument is \texttt{P02.1-argument.md}. The saved small
program \texttt{programs/check\_finite\_net.py} exhaustively checks PL
monotonicity in a finite five-branch proxy using exact rational arithmetic;
it is not an analytic or uncountability verification. Saved rerun commands
and environment details are in \texttt{rerun.sh} and \texttt{evidence/}.
Compilation and execution success do not independently verify the proof.
A separate source-aware audit using \texttt{gpt-6.1-sol} at
\texttt{xhigh} returned PASS for this exact partial result; the report
and premises are retained in \texttt{P02.1-review.json}.

\begin{thebibliography}{9}
\bibitem{GR} W. Gubler and J. Rabinoff,
\emph{Classically psh and pluriharmonic functions on Berkovich spaces},
arXiv:2506.13548v2, 17 September 2025.
\bibitem{Ber} V. G. Berkovich,
\emph{Spectral theory and analytic geometry over non-Archimedean fields},
Mathematical Surveys and Monographs 33, American Mathematical Society,
1990.
\end{thebibliography}
\end{document}
